\section*{Supplementary material}
Supplementary S1 records protocol changes; S2 gate definitions and operating curves; S3 dataset geometry; S4 result-to-artifact traceability; S5 slip analysis; and S6 the implementation audit, nominal statistical tests and sensitivity analyses. A reproducibility package with the scripts and result summaries behind S6 accompanies the submission.

\section*{CRediT authorship contribution statement}
\textbf{Mohammad Rafiur Rahman:} Conceptualization, Methodology, Software, Validation, Formal analysis, Investigation,
Data curation, Writing -- original draft, Writing -- review \& editing, Visualization, Project administration.

\section*{Declaration of competing interest}

The author declares no known competing financial interests or personal relationships that could have
appeared to influence the work reported in this paper.

\section*{Funding}
This research did not receive any specific grant from funding agencies in the public, commercial, or not-for-profit
sectors.

\section*{Data availability}

All raw data are public. Paderborn bearing data are distributed by the KAt-DataCenter, Universit\"at Paderborn
\cite{lessmeier2016}; HUST bearing is available from Mendeley Data (\url{https://doi.org/10.17632/cbv7jyx4p9})
\cite{thuan2023hust}; CWRU bearing data from the Case Western Reserve University Bearing Data Center \cite{cwrubdc};
UORED-VAFCLS from the University of Ottawa \cite{sehri2023uored}; and the JNU data are the Jiangnan University bearing data
commonly cited to Li et al.\ \cite{li2013jnu}.
The research data produced in this study are available at \url{https://github.com/rafilovestosuffer/physgate-sfda} and archived
on Zenodo as the experimental release v1.0-submission, \url{https://doi.org/10.5281/zenodo.23105810} \cite{rahman2026physgate}.
They comprise the frozen split definitions, the per-run experiment ledger (one row per run and task, each with its kernel,
commit and artifact path), the analysis scripts with their checksums, the version-controlled protocol and the retained run
artifacts that back every number in this paper. The audit scripts and result summaries behind Supplementary~S6 were prepared after that release and are supplied with this submission.

\section*{Code availability}

The evaluation protocol, gate implementations, analysis and figure scripts are released with the research data. SDALR is
run from its authors' released code (BdLab405/SDALR, commit fb9c379), which the repository fetches by script because
it carries no licence for redistribution; the evaluation wrapper and the gate hook described in Section~\ref{sec:protocol} are
added around it without changing its training code. SHOT runs from our implementation of its objective on the same backbone (runners/shot\_patch.py, version 2); the first version, whose indexing defect is documented in Supplementary~S6, is retained for audit only and its outputs are not used.

\section*{Declaration of generative AI and AI-assisted technologies in the manuscript preparation process}

During the preparation of this work the author used Claude (Anthropic; through the Claude Code agent, with the model
claude-opus-5-5) in order to draft and edit text, write analysis and figure code, draw the schematic
figures with plotting code, and check numerical consistency and references against the retained artifacts and primary
sources. OpenAI Codex was used to audit the implementation, reaggregate retained results, revise the text and rebuild the submission files; it ran no training or experiment. After using these tools, the author reviewed and edited the content as needed and takes full responsibility for the
content of the publication. Because the Claude workflow also executed analyses and compute jobs, that use forms part of the
research method and is described in Section~\ref{sec:protocol}.
